\documentclass[tikz,border=10pt]{standalone}
\usetikzlibrary{positioning, arrows.meta, calc, babel, fit}
\usepackage[utf8]{inputenc}
\usepackage[T1]{fontenc}
\usepackage{lmodern}
\usepackage[spanish,es-tabla]{babel}
\usepackage{amsmath, amssymb, bm}

\begin{document}
\begin{tikzpicture}[
  node distance=0.6cm and 0.8cm,
  align=center,
  font=\sffamily\small,
  % Estilos de bloques principales
  box/.style={draw, rectangle, rounded corners=2pt, minimum width=2.2cm, minimum height=0.9cm, thick, fill=white},
  mlpbox/.style={draw=blue!80!black, rectangle, rounded corners=2pt, minimum width=3.2cm, minimum height=1.0cm, thick, fill=blue!5},
  rnnbox/.style={draw=teal!80!black, rectangle, rounded corners=2pt, minimum size=0.8cm, thick, fill=teal!10},
  cellbox/.style={draw=purple!80!black, rectangle, rounded corners=6pt, dashed, thick, fill=purple!2},
  grubox/.style={draw=orange!80!black, rectangle, rounded corners=6pt, dashed, thick, fill=orange!2},
  % Operaciones matemáticas internas
  op/.style={draw, circle, minimum size=0.45cm, inner sep=0pt, thick, fill=white, font=\sffamily\scriptsize},
  gate/.style={draw, rectangle, rounded corners=1pt, minimum width=1.25cm, minimum height=0.65cm, fill=gray!10, font=\sffamily\scriptsize},
  % Estilos de flechas
  line/.style={draw, -{Stealth[scale=1.2]}, thick},
  gradline/.style={draw, {Stealth[scale=1.2]}-, dashed, red!80!black, thick},
  labeltext/.style={font=\sffamily\small, text=black!80}
]

  % ==========================================
  % PANEL A: CASO FEEDFORWARD (MLP) - SIN MEMORIA
  % ==========================================
  \node[font=\sffamily\bfseries\large, text=black!85, anchor=center] (tituloA) at (11.2, 8.5) {A. Perceptrón Multicapa (MLP) -- Sin memoria temporal};
  
  \node[mlpbox] (nnA) at (11.2, 7.3) {Red MLP\\(Capas Lineales + ReLU)};
  \node[box, draw=blue!50!black, fill=gray!3, left=1.8cm of nnA] (inputA) {Entrada actual\\$\bm{x}_k \in \mathbb{R}^9$};
  \node[box, draw=gray!50!black, fill=gray!3, right=1.8cm of nnA] (outputA) {Fuerza deseada\\$\widehat{\bm{F}}_{d,W,k} \in \mathbb{R}^3$};

  \draw[line] (inputA) -- (nnA);
  \draw[line] (nnA) -- (outputA);

  % ==========================================
  % PANEL B: RED RECURRENTE BÁSICA (RNN) - MEMORIA SIMPLE
  % ==========================================
  \node[font=\sffamily\bfseries\large, text=black!85, anchor=center] (tituloB) at (11.2, 5.0) {B. Red Recurrente Básica (RNN) -- Memoria lineal simple sin compuertas};

  % Desenrollado temporal k-2, k-1, k
  \node[rnnbox] (h0) at (6.95, 3.6) {Estado\\$h_{k-2}$};
  \node[rnnbox, right=2.0cm of h0] (h1) {Estado\\$h_{k-1}$};
  \node[rnnbox, right=2.0cm of h1] (h2) {Estado\\$h_{k}$};
  
  \node[box, draw=gray!50, fill=gray!1, minimum width=1.5cm, minimum height=0.6cm, below=0.6cm of h0] (x0) {$\bm{x}_{k-2}$};
  \node[box, draw=gray!50, fill=gray!1, minimum width=1.5cm, minimum height=0.6cm, below=0.6cm of h1] (x1) {$\bm{x}_{k-1}$};
  \node[box, draw=gray!50, fill=gray!1, minimum width=1.5cm, minimum height=0.6cm, below=0.6cm of h2] (x2) {$\bm{x}_{k}$};

  \node[box, draw=gray!50!black, fill=gray!3, right=1.5cm of h2] (outputB) {Fuerza deseada\\$\widehat{\bm{F}}_{d,W,k}$};

  % Conexiones verticales
  \draw[line] (x0) -- (h0);
  \draw[line] (x1) -- (h1);
  \draw[line] (x2) -- (h2);

  % Conexiones horizontales (paso del tiempo)
  \draw[line] ($(h0.east) + (0,0.15)$) -- node[above=0.02cm, font=\sffamily\scriptsize, text=black!60] {$h_{t-1} \to h_t$} ($(h1.west) + (0,0.15)$);
  \draw[line] ($(h1.east) + (0,0.15)$) -- ($(h2.west) + (0,0.15)$);
  \draw[line] (h2) -- (outputB);
  
  % Conexiones del gradiente (retropropagación hacia atrás)
  \draw[gradline, red!60!black, line width=1.2pt] ($(h1.west) + (0,-0.15)$) -- node[below, font=\sffamily\tiny\bfseries, text=red!80!black] {Gradiente} ($(h0.east) + (0,-0.15)$);
  \draw[gradline, red!30!black, line width=0.6pt] ($(h2.west) + (0,-0.15)$) -- node[below, font=\sffamily\tiny, text=red!60!black] {Atenuación} ($(h1.east) + (0,-0.15)$);
  
  % Entrada y salida del flujo temporal general
  \draw[line, dashed, gray] ($(h0.west) + (-0.8,0.15)$) -- ($(h0.west) + (0,0.15)$);
  \node[left=0.8cm of h0, yshift=0.15cm, font=\sffamily\scriptsize, text=black!60] {$h_{k-3}$};

  % ==========================================
  % PANEL C: ARQUITECTURAS CON COMPUERTAS (LSTM / GRU) - MEMORIA SELECTIVA
  % ==========================================
  \node[font=\sffamily\bfseries\large, text=black!85, anchor=center] (tituloC) at (11.2, 0.2) {C. Arquitecturas Recurrentes con Compuertas -- Regulación selectiva de la memoria};

  % --- SUB-PANEL LSTM ---
  \coordinate (lstmorigen) at (0.2, -1.0);

  \node[cellbox, minimum width=10.2cm, minimum height=6.0cm, anchor=north west] (lstmcella) at (lstmorigen) {};
  \node[above=0.05cm of lstmcella.north west, anchor=south west, font=\sffamily\bfseries\small, text=purple!80!black] {Célula LSTM};

  % Carril superior: estado de celda
  \node[labeltext, anchor=east] (c_prev) at ($(lstmorigen) + (0, -1.0)$) {$c_{k-1}$};
  \node[op] (mult_f) at ($(lstmorigen) + (2.6, -1.0)$) {$\otimes$};
  \node[op] (add_i) at ($(lstmorigen) + (5.8, -1.0)$) {$\oplus$};
  \node[labeltext, anchor=west] (c_next) at ($(lstmorigen) + (10.2, -1.0)$) {$c_k$};

  % Fila de compuertas (columnas alineadas con operaciones)
  \node[gate] (gate_f) at ($(lstmorigen) + (2.6, -3.3)$) {Compuerta\\olvido ($f_k$)};
  \node[gate] (gate_i) at ($(lstmorigen) + (5.8, -3.3)$) {Compuerta\\entrada ($i_k$)};
  \node[gate] (gate_o) at ($(lstmorigen) + (7.8, -3.3)$) {Compuerta\\salida ($o_k$)};

  % Carril de salida oculta (columna derecha, por encima de compuertas)
  \node[labeltext, anchor=east] (h_prev) at ($(lstmorigen) + (0, -2.1)$) {$h_{k-1}$};
  \node[op] (tanh_act) at ($(lstmorigen) + (7.6, -2.1)$) {$\tanh$};
  \node[op] (mult_o) at ($(lstmorigen) + (9.3, -2.1)$) {$\otimes$};
  \node[labeltext, anchor=west] (h_next) at ($(lstmorigen) + (10.2, -2.1)$) {$h_k$};

  \node[box, draw=gray!40, fill=gray!1, minimum width=1.0cm, minimum height=0.5cm] (x_lstm) at ($(lstmorigen) + (5.0, -6.35)$) {$\bm{x}_k$};

  % Flujo de estado de celda (carril alto, sin cruces)
  \draw[line] (c_prev.east) -- (mult_f.west);
  \draw[line] (mult_f.east) -- (add_i.west);
  \draw[line] (add_i.east) -- (c_next.west);

  % Compuertas al carril superior (verticales cortas)
  \draw[line] (gate_f.north) -- (mult_f.south);
  \draw[line] (gate_i.north) -- (add_i.south);

  % Salida oculta: rama de $c_k$ entra a tanh por la izquierda
  \draw[line] (add_i.east) -- ++(0.5, 0) coordinate (lstm_branch_c);
  \draw[line] (lstm_branch_c) |- (tanh_act.west);
  \draw[line] (tanh_act.east) -- (mult_o.west);
  \draw[line] (mult_o.east) -- (h_next.west);

  % Compuerta de salida: sale por el este, baja y entra al $\otimes$ desde abajo
  \draw[line] (gate_o.east) -- ++(0.55, 0) coordinate (lstm_o_east);
  \draw[line] (lstm_o_east) -| (mult_o.south);

  % $h_{k-1}$: anillo inferior amplio; entra a compuertas por la derecha
  \coordinate (lstm_hdrop) at ($(lstmorigen) + (1.1, -2.1)$);
  \coordinate (lstm_hwrap) at ($(lstmorigen) + (1.1, -6.7)$);
  \coordinate (lstm_hring) at ($(lstmorigen) + (8.53, -6.7)$);
  \draw[line] (h_prev.east) -- (lstm_hdrop);
  \draw[line] (lstm_hdrop) -- (lstm_hwrap);
  \draw[line] (lstm_hwrap) -- (lstm_hring);
  \draw[line] (lstm_hwrap) -| (gate_f.south east);
  \draw[line] (lstm_hwrap) -| (gate_i.south east);
  \draw[line] (lstm_hring) -| (gate_o.south east);

  % $x_k$: bus intermedio; entra a compuertas por la izquierda
  \coordinate (lstm_xbus) at ($(x_lstm.north) + (0, 0.45)$);
  \draw[line] (x_lstm.north) -- (lstm_xbus);
  \draw[line] (lstm_xbus) -| (gate_f.south west);
  \draw[line] (lstm_xbus) -| (gate_i.south west);
  \draw[line] (lstm_xbus) -| (gate_o.south west);


  % --- SUB-PANEL GRU ---
  \coordinate (gruorigen) at (13.0, -1.0);

  \node[grubox, minimum width=9.2cm, minimum height=6.0cm, anchor=north west] (grucella) at (gruorigen) {};
  \node[above=0.05cm of grucella.north west, anchor=south west, font=\sffamily\bfseries\small, text=orange!80!black] {Célula GRU};

  % Carril superior: actualización de $h_k$
  \node[labeltext, anchor=east] (h_gru_prev) at ($(gruorigen) + (0, -1.2)$) {$h_{k-1}$};
  \node[op] (mix) at ($(gruorigen) + (3.2, -1.2)$) {$\otimes$};
  \node[op] (add_gru) at ($(gruorigen) + (6.8, -1.2)$) {$\oplus$};
  \node[labeltext, anchor=west] (h_gru_next) at ($(gruorigen) + (9.2, -1.2)$) {$h_k$};

  % Fila inferior: compuertas y candidato (más separados)
  \node[gate] (gate_z) at ($(gruorigen) + (2.2, -3.4)$) {Compuerta\\act. ($z_k$)};
  \node[gate] (gate_r) at ($(gruorigen) + (4.8, -3.4)$) {Compuerta\\reinicio ($r_k$)};
  \node[box, draw=orange!40, fill=orange!2, minimum width=1.6cm, minimum height=0.6cm, font=\sffamily\scriptsize] (cand) at ($(gruorigen) + (7.8, -3.4)$) {Candidato $\tilde{h}_k$};

  \node[box, draw=gray!40, fill=gray!1, minimum width=1.0cm, minimum height=0.5cm] (x_gru) at ($(gruorigen) + (4.5, -6.35)$) {$\bm{x}_k$};

  % Flujo principal de $h_k$ (margen izquierdo, por encima de compuertas)
  \coordinate (gru_hdrop) at ($(gruorigen) + (1.1, -1.2)$);
  \draw[line] (h_gru_prev.east) -- (gru_hdrop);
  \draw[line] (gru_hdrop) -- (mix.west);
  \draw[line] (mix.east) -- (add_gru.west);
  \draw[line] (add_gru.east) -- (h_gru_next.west);

  % Compuerta de actualización al $\otimes$ con codo ortogonal
  \draw[line] (gate_z.north) -- ++(0, 0.55) coordinate (gru_z_elbow);
  \draw[line] (gru_z_elbow) -| (mix.south);

  % Candidato al sumador (columna propia, separada del reinicio)
  \draw[line] (gate_r.east) -- (cand.west);
  \draw[line] (cand.north) -- ++(0, 0.575) -| (add_gru.south);

  % $h_{k-1}$: anillo inferior amplio; entra por la derecha de compuertas y candidato
  \coordinate (gru_hwrap) at ($(gruorigen) + (1.1, -6.7)$);
  \coordinate (gru_hring) at ($(gruorigen) + (8.5, -6.7)$);
  \draw[line] (gru_hdrop) -- (gru_hwrap);
  \draw[line] (gru_hwrap) -- (gru_hring);
  \draw[line] (gru_hwrap) -| (gate_z.south east);
  \draw[line] (gru_hwrap) -| (gate_r.south east);
  \draw[line] (gru_hring) -| (cand.south east);

  % $x_k$: bus intermedio; entra por la izquierda
  \coordinate (gru_xbus) at ($(x_gru.north) + (0, 0.45)$);
  \draw[line] (x_gru.north) -- (gru_xbus);
  \draw[line] (gru_xbus) -| (gate_z.south west);
  \draw[line] (gru_xbus) -| (gate_r.south west);
  \draw[line] (gru_xbus) -| (cand.south west);

\end{tikzpicture}
\end{document}
